\chapter{Continuous and Runtime-Managed Compilation}
\label{ch:continuous-compilation}

A runtime can finish compiling a function just after the source or an assumed class layout changes.
Publishing that finished artifact because the worker succeeded would install stale code.
Rejecting the stale artifact solves only half the problem: an older version may still be executing and cannot yet be reclaimed.
Continuous compilation therefore needs a lifecycle for artifacts, not just a queue of optimization jobs.

We use \term{continuous compilation} for an architecture in which compilation, execution, observation, invalidation, and replacement overlap throughout a process lifetime.
This is a design description rather than an attribution to an unidentified framework.
Established components include incremental queries, tiered execution, lazy materialization, deoptimization, and versioned publication.
The chapter combines them into a small explicit ownership model.

\section{Model several representations without inventing a lattice}
\label{sec:ch67-representations}

A function can have source syntax, a typed graph, bytecode, baseline code, specialized code, and target-specific variants.
These representations do not automatically form a mathematical lattice.
To call them a lattice requires a partial order and a least upper bound and greatest lower bound for every pair.
Optimization level alone provides none of these guarantees: one version may use less memory while another runs faster on a different device.

An \term{artifact graph} is the more direct model.
Nodes carry representation kind, source revision, target, assumptions, entry contract, and ownership information.
Edges describe derivation, dependencies, dispatch, or valid state transfers, with edge kinds kept distinct.
A derivation edge does not imply that a running frame can jump between the two nodes.
Figure~\ref{fig:ch67-artifacts} shows a branching representation graph.

\begin{figure}[htbp]
\centering
\begin{tikzpicture}[node distance=16mm and 18mm]
\node[basic block] (ir) {semantic IR, revision r};
\node[basic block,below left=of ir] (byte) {portable bytecode};
\node[basic block,below right=of ir] (native) {CPU specialization};
\node[basic block,below=of byte] (interp) {interpreter entry};
\node[basic block,below=of native] (gpu) {device variant};
\draw[cfg edge] (ir)--(byte); \draw[cfg edge] (ir)--(native);
\draw[cfg edge] (byte)--(interp);
\draw[cfg edge] (ir.east) -| ([xshift=12mm]native.east) |- (gpu.east);
\end{tikzpicture}
\caption{Representation choices branch by target and contract; they need not have a single monotone optimization order.}
\label{fig:ch67-artifacts}
\end{figure}

An execution graph can overlay observations such as invocation counts, transitions, and version use.
Those observations guide policy but do not replace semantic dependencies.
A hot edge is evidence about frequency; a call edge or invalidation dependency is a correctness-relevant relation.
Conflating them would make an unobserved path disappear from the compiler's meaning of the program.

\section{Publish only against a matching epoch}
\label{sec:ch67-publication}

The companion runtime has one serialized owner and an integer epoch.
A worker conceptually receives an immutable request labeled with the current epoch.
An edit advances the epoch and installs a generic version for new entries.
When a worker returns, publication compares its request epoch with the owner's current epoch.
Algorithm~\ref{alg:ch67-publish} rejects a mismatch.

\begin{algorithm}[htbp]
\caption{Publish a completed version through one serialized owner.}
\label{alg:ch67-publish}
\begin{pseudocode}
\Require validated artifact $A$ with request epoch $e$, current epoch $r$
\Ensure stale results never replace the current entry
\State enter the owner operation
\If{$e\ne r$}
  \State discard the unpublished artifact; \Return stale
\EndIf
\State fully initialize entry metadata and recovery dependencies for $A$
\State mark the previous entry retired
\State install $A$ as the current entry
\State \Return published
\end{pseudocode}
\end{algorithm}

Epoch comparison is constant-time in the model.
Artifact validation and initialization have their own costs and happen before publication becomes visible.
A production request key may include separate source, target, layout, and configuration versions rather than one global epoch.
The global epoch is simple but invalidates more work than necessary.

Table~\ref{tab:ch67-epoch} works a stale completion.
An old active invocation continues under its old-version contract in this model; new entries select the new generic version.
This is a defined edit policy, not a claim that arbitrary language-level assumptions remain valid for active optimized frames.
If an assumption affects an existing frame's safety, the runtime must arrange a safepoint, patch, guard failure, or other supported deoptimization protocol before the assumption becomes false.

\begin{table}[htbp]
\caption{A worker's successful computation can still be stale at publication time.}
\label{tab:ch67-epoch}
\centering
\begin{tabular}{lll}
\toprule
Event & Owner epoch & Publication consequence \\
\midrule
worker starts & $0$ & request carries $0$ \\
input edit & $1$ & generic version becomes current \\
old worker returns & $1$ & reject request epoch $0$ \\
new worker returns & $1$ & epoch $1$ may publish \\
\bottomrule
\end{tabular}
\end{table}

\section{Separate retirement from reclamation}
\label{sec:ch67-reclamation}

\term{Retirement} prevents new entries from selecting an artifact.
\term{Artifact reclamation} releases its code and metadata once no permitted use remains.
An active frame, a return address, an inline cache, a pending compilation, or another version's recovery map may keep an artifact reachable.
Freeing code because it is no longer the preferred entry is therefore incorrect.

The teaching runtime grants a lease when an invocation acquires the current version.
Acquisition and retirement run through the same serialized owner, eliminating the race between reading an entry and incrementing its use count.
A retired version is collected only after all leases have been released.
Algorithm~\ref{alg:ch67-collect} states the modeled rule.

\begin{algorithm}[htbp]
\caption{Reclaim retired artifacts in a lease-complete serialized model.}
\label{alg:ch67-collect}
\begin{pseudocode}
\Require every possible use owns a lease, version list, serialized owner
\Ensure no artifact with a possible outstanding use is reclaimed
\ForAll{versions $v$}
  \If{$v$ is retired and its lease count is zero}
    \State remove $v$ from the owned version list
    \State release its code and metadata resources
  \EndIf
\EndFor
\end{pseudocode}
\end{algorithm}

Collection scans $V$ versions in $O(V)$ time, excluding resource destruction.
The Python implementation models active invocation leases only; it contains symbolic code bodies rather than executable allocations.
The stronger pseudocode precondition identifies what must be added when code pointers can escape.
Reference cycles among artifacts also require a cycle-aware ownership strategy or tracing; active-frame counters alone do not discover them.

Figure~\ref{fig:ch67-lifetime} distinguishes the lifetime states.
An epoch-based reclamation scheme, hazard pointers, or stop-the-world scanning can implement a different mechanism, but each must close the same acquisition/reclamation race.
Calling a memory-management technique by name does not prove that all code references are covered.

\begin{figure}[htbp]
\centering
\begin{tikzpicture}[node distance=14mm]
\node[basic block] (new) {validated, unpublished};
\node[basic block,below=of new] (current) {current entry};
\node[basic block,below=of current] (retired) {retired, possibly leased};
\node[basic block,below=of retired] (free) {reclaimed};
\draw[cfg edge] (new)--(current); \draw[cfg edge] (current)--(retired);
\draw[cfg edge] (retired)--node[right,font=\small] {no remaining use}(free);
\end{tikzpicture}
\caption{Replacing an entry retires its old artifact before reclamation becomes possible.}
\label{fig:ch67-lifetime}
\end{figure}

\section{Migrate data behind a stable handle}
\label{sec:ch67-migration}

Representation evolution also applies to data.
A mostly zero array might move from dense storage to a sparse map while preserving its logical length and element values.
If clients hold a stable handle and access through it, the runtime can build a replacement privately and switch the handle after completion.
The companion retains array identity while alternating between a list and a map that omits zero entries.

For logical values $[0,4,0,1,0]$, the sparse payload is $\{1\mapsto4,3\mapsto1\}$ and the length remains five.
An alias of the handle still reads four at index one and zero at index four.
Index five must still fail; absence from the sparse map does not mean every integer is a valid array index.
Table~\ref{tab:ch67-layout} records the distinction between identity and payload.

\begin{table}[htbp]
\caption{Migration changes physical storage while preserving the logical array contract.}
\label{tab:ch67-layout}
\centering
\begin{tabular}{lll}
\toprule
Property & Dense form & Sparse form \\
\midrule
Logical length & $5$ & $5$ \\
Stored payload values & $0,4,0,1,0$ & indices $1,3$ hold $4,1$ \\
Read at index $4$ & $0$ & $0$ \\
Read at index $5$ & error & error \\
Handle identity & preserved & preserved \\
\bottomrule
\end{tabular}
\end{table}

The model is sequential and the switch consists of ordinary host assignments.
A concurrent implementation needs atomic publication of a coherent descriptor, exclusion or reconciliation of concurrent writes, and a lifetime rule for readers of the old payload.
Native code that bypasses the handle also needs a layout guard or invalidation dependency.
Thus data migration and code invalidation are coupled even when the logical value is unchanged.

\section{Coordinate actors and compilation services}
\label{sec:ch67-actors}

An actor-style architecture can give one owner responsibility for each mutable compilation unit.
Workers receive immutable snapshots and return proposals; the owner validates and publishes them.
This avoids shared mutable IR inside workers, but messages can arrive out of order and duplicate work can finish twice.
Request IDs, epochs, idempotent completion handling, and cancellation semantics remain necessary.

A persistent compilation service may reuse parsed modules, summaries, profiles, and generated artifacts across requests.
Its cache key must describe the semantic and target configuration completely.
Resource budgets need to cover retained IR, code versions, pending jobs, and metadata, not just emitted machine-code bytes.
Backpressure can delay low-value jobs; cancellation should discard unpublished work without abandoning resources already reachable by execution.

The runtime owns optimization policy when it chooses what to compile, where to execute it, and when to retire it using live observations.
The compiler still owns legality checks, and the publication owner still owns consistency.
A profitability estimate can decide to keep generic code indefinitely.
It cannot justify installing an artifact with an invalid assumption or reclaiming an active frame's continuation.

\section{Exercises}
\label{sec:ch67-exercises}

\begin{enumerate}
\item\label{ex:ch67-stale} Acquire epoch-zero code, start a compilation, edit to epoch one, and complete the old request.
State the publication result and when the acquired version can be reclaimed.
\item\label{ex:ch67-race} Show the race if a thread reads a code pointer before incrementing its lease while another thread retires and collects it.
Explain how serialized acquisition closes the race.
\item\label{ex:ch67-sparse} Migrate $[0,4,0,1,0]$ to sparse storage and back.
Give the payload, valid index range, and two alias observations that must remain unchanged.
\item\label{ex:ch67-lattice} Two versions trade runtime speed against code size.
Explain why this does not define a representation lattice and name a relation the artifact graph can record without that assumption.
\end{enumerate}

\section{Further reading}
\label{sec:ch67-reading}

ORC documents modular JIT materialization and resource-management interfaces~\cite{llvm-2026-orc}.
The incremental-query and deoptimization references supply the dependency and state-transfer components of this architecture~\cite{rust-2026-incremental,fluckiger-2018-deoptimization}.
The companion's lease and migration models test these small contracts without claiming to implement a concurrent native JIT service.
